\documentclass[11pt,a4paper]{article}
\usepackage[spanish]{babel}
\usepackage{fontspec}
\setmainfont[Path=../assets/fonts/,UprightFont=InterVariable.ttf,BoldFont=InterVariable.ttf,ItalicFont=InterVariable.ttf,BoldItalicFont=InterVariable.ttf]{InterVariable.ttf}
\setsansfont[Path=../assets/fonts/]{InterVariable.ttf}
\usepackage[margin=25mm]{geometry}
\usepackage{amsmath,amssymb,booktabs}
\usepackage{xcolor,hyperref}
\definecolor{usmblue}{HTML}{003876}
\hypersetup{colorlinks=true,linkcolor=usmblue,urlcolor=usmblue}
\newcommand{\E}{\mathbb E}
\newcommand{\Prob}{\mathbb P}
\title{\textcolor{usmblue}{Control predictivo estocástico informado por riesgo}\\
\large Estado técnico y primer hito de certificación de costos}
\author{Proyecto DT--SMDFR Risk Control\\Universidad Técnica Federico Santa María}
\date{22 de octubre de 2026}
\begin{document}
\maketitle
\begin{abstract}
Se audita la distribución v1.0.0 de leomiq/dt\_riskcontrol\_smdfr,
commit 3266149. El controlador minimiza un costo de seguimiento y márgenes
deterministas mediante Luus--Jaakola o algoritmo genético; el margen
probabilístico por IS--MCDET se obtiene en posprocesamiento. Se formula
una cota inferior por intervalos para el costo vigente y una ruta hacia
optimización estocástica causal con varios sistemas de protección.
\end{abstract}

\section{Estado y brechas del código}
El piloto \texttt{pilot\_single\_run.m} llama
\texttt{luus\_jaakola\_det.m}, \texttt{DET\_2.m} y
\texttt{fitness\_cost.m} para cada candidato. \texttt{ga\_det.m}
emplea la misma función de costo. En
\nolinkurl{DET_scheduler_v1_perc_online.m} las llamadas LJ/GA están
comentadas y se usa PID; por ello el piloto es la referencia de búsqueda.
El README registra $f=0.326151$ a $t_q=1180$ s frente a $0.326139$
en el árbol de referencia, diferencia relativa cercana a $4\,10^{-5}$.

La función de costo concatena muestras futuras de nodos de
\texttt{forest\_data\{1\}} y evalúa
\begin{equation}
J(u)=
\underbrace{\sum_{j\in\mathcal I(u)}(T_j(u)-T_{\rm ref})^2}_{J_T}
+\underbrace{\lambda_m\sum_{k=1}^{K-1}
 [\,|m_{k+1}(u)-m_k(u)|-\Delta_m\,]_+^2}_{J_m}
+\underbrace{\lambda_s\sum_{j\in\mathcal I(u)}
 [\,g(T_j(u))\,]_+^2}_{J_s}.
\label{eq:old}
\end{equation}
Aquí $m_k$ procede de \texttt{CurrentToMassFlow.m},
$\Delta_m=5$, $\lambda_m=10^{-4}$ y $[x]_+=\max(x,0)$.
Por defecto $g(T)=T-U_y$, $U_y=1244$ K y $\lambda_s=10$.
Con \texttt{dsm\_thr} y \texttt{lambda\_dsm} en el nodo raíz se usa
$g(T)=(T-T_{\rm ref})/(U_y-T_{\rm ref})-\tau$.
El booleano \texttt{feasible} comprueba
$|(T-T_{\rm ref})/(U_y-T_{\rm ref})|\leq1$ en las muestras,
pero LJ registra candidatos no factibles y GA sí los filtra.
La semántica de factibilidad debe unificarse antes de certificar el
óptimo restringido.

La ecuación \eqref{eq:old} no es una esperanza: concatena nodos de un
solo árbol y omite las probabilidades de ruta. Incluso el número de
muestras puede variar con la topología inducida por el control.
El PSM se calcula posteriormente en
\nolinkurl{DET_scheduler_postprocessing_perc_online.m} mediante
trayectorias terminales y pesos de importancia. El DSM penalizado dentro
del costo es determinista; no equivale a una métrica probabilística
evaluada para cada candidato.

\section{Cota inferior del costo vigente}
Sea $\mathcal U_b$ una celda de controles. Para toda $u\in\mathcal U_b$,
supónganse encierros rigurosos
$T_j(u)\in[L_j,H_j]$ para cada muestra futura de una topología DET fija,
y $m_k(u)\in[\ell_k,h_k]$. Definamos
\[
\begin{aligned}
d_j&=\max\{L_j-T_{\rm ref},T_{\rm ref}-H_j,0\},\\
a_k&=\ell_{k+1}-h_k,\qquad b_k=h_{k+1}-\ell_k,\\
e_k&=\max\{a_k-\Delta_m,-\Delta_m-b_k,0\}.
\end{aligned}
\]
Como la penalización de seguridad crece con $T$, su mínimo por intervalo
se alcanza en $L_j$. Así,
\begin{equation}
\begin{aligned}
\underline J(\mathcal U_b)
&=\sum_j d_j^2+\lambda_m\sum_k e_k^2\\
&\quad+
\begin{cases}
\lambda_s\sum_j[L_j-U_y]_+^2,&\text{límite absoluto},\\
\lambda_s\sum_j
 [\,\frac{L_j-T_{\rm ref}}{U_y-T_{\rm ref}}-\tau\,]_+^2,
 &\text{penalización DSM}
\end{cases}\\
&\leq\inf_{u\in\mathcal U_b}J(u).
\end{aligned}
\label{eq:lb}
\end{equation}
La cota puede ser estricta: los mínimos individuales pueden requerir
controles diferentes. \nolinkurl{certified_cost_lower_bound_box.m}
implementa \eqref{eq:lb}. En cajas puntuales reproduce exactamente
\eqref{eq:old}, salvo redondeo.

\paragraph{Alcance.}
El módulo no construye los encierros de temperatura. Una malla,
semillas repetidas o un surrogate no los certifican. Si la topología
cambia, hay que cubrir todas sus alternativas con una convención de
horizonte común. Si se obtiene un candidato factible con costo $U$ y
una partición exhaustiva con cotas $L_b$, entonces
$L=\min_bL_b\leq J^\star\leq U$ y $U-L$ certifica el gap.
Sin encierros válidos, sólo se tiene $J\geq0$.
La prueba sintética \nolinkurl{test_certified_cost_lower_bound_box.m}
verifica exactitud puntual y validez de inclusión.

\section{Problema estocástico propuesto}
Sea $\pi$ una política causal; el estado híbrido $z_k$ contiene modos
de cada protección y del sistema pasivo. Los eventos $\xi_k$ incluyen
fallo, reparación, perturbaciones e incertidumbre del modelo.
Una formulación inicial es
\begin{align}
\min_{\pi}\quad&
\E_p\!\left[\sum_{k=0}^{N-1}
 w_T(T_k-T_{\rm ref})^2+w_u c_u(u_k,a_k)\right]
+\lambda_R R(\pi),\\
\text{sujeto a}\quad&
\Prob_p\!\left(\max_{0\leq k\leq N}T_k>U_y
\mid\mathcal F_t\right)\leq\varepsilon,\\
&x_{k+1}=F(x_k,z_k,u_k,a_k,\xi_k),\quad
z_{k+1}=G(z_k,a_k,\xi_k).
\end{align}
Debe imponerse no anticipación en ramas aún indistinguibles.
Son candidatos para $R$ la probabilidad de excedencia, CVaR del exceso
térmico y un PSM basado en cuantiles; no son intercambiables sin fijar
horizonte, confianza y unidades. En muestreo por importancia,
$w_i=p(\omega_i)/q(\omega_i)$ debe representar la trayectoria completa,
incluidas probabilidades condicionadas a demanda y acción. Comparar
políticas requiere tamaño efectivo de muestra, intervalos de confianza
y evaluación fuera de muestra.

\section{Protecciones múltiples y ruta de implementación}
El árbol actual contiene el par \texttt{[demand\_ps,ava\_ps]} del ACS.
Con un solo actuador activo no surge una decisión de orden entre
protecciones. Al incorporar otros sistemas, $z_k$ debe incluir
disponibilidad, demanda, activación, fallo, reparación y retardos;
el sistema pasivo tiene demanda y éxito propios. La secuencia temporal
modifica acciones factibles, respuesta térmica y probabilidades.
Deben representarse causas comunes y recursos compartidos.

\begin{center}
\begin{tabular}{@{}p{12mm}p{68mm}p{63mm}@{}}
\toprule
Hito & Entrega & Cierre verificable\\
\midrule
M0 & Piloto y desglose de costo & SHA, semilla y diferencia frente al valor archivado\\
M1 & Encierros ROM/DET y cotas por celdas & $L\leq J^\star\leq U$ y gap, o limitación explícita\\
M2 & Riesgo dentro de cada candidato & Pesos IS validados y decisiones sensibles al riesgo\\
M3 & Comparación estocástica & Frontera costo--riesgo fuera de muestra y latencia\\
M4 & Múltiples sistemas y pasivo & Secuencias, fallos comunes y política causal trazable\\
\bottomrule
\end{tabular}
\end{center}
El desarrollo y los criterios de aceptación constan en
\texttt{docs/implementation\_route.md}.
\end{document}
